% \documentclass[../main]{subfiles}
\ifSubfilesClassLoaded{\setcounter{chapter}{9}}{}
% \begin{document}

\chapter{Pembangkit Bilangan Acak dan Steganografi}

\begin{subcpmk}
  \item Sub-CPMK 6.1: Membedakan PRNG dan TRNG serta menguji keacakan deret bit
  \item Sub-CPMK 7.1: Mengimplementasikan metode LSB untuk steganografi pada citra digital
\end{subcpmk}

% ============================================================
% MATERI POKOK
% ============================================================
\section{Asesmen Komprehensif}

Bab ini berisi asesmen untuk mengukur pencapaian seluruh CPMK yang telah dipelajari sepanjang mata kuliah Kriptografi. Menurut sumber evaluasi, asesmen komprehensif dirancang untuk mengukur pencapaian mahasiswa secara holistik, mencakup aspek kognitif, psikomotorik, dan afektif yang terukur dalam pembelajaran kriptografi.

\subsection{Soal UTS (Style)}

Soal UTS dirancang untuk mengukur tiga pilar utama pencapaian CPMK. Menurut sumber evaluasi, setiap soal harus dihubungkan dengan Sub-CPMK secara eksplisit dan diukur dengan kriteria penilaian yang jelas. Soal UTS yang baik mengukur kemampuan analitis, sintesis, dan evaluasi konsep kriptografi secara terintegrasi.

\begin{enumerate}
  \item \textbf{Jelaskan tiga pilar keamanan informasi} (confidentiality, integrity, authenticity) dan berikan contoh ancaman yang melanggar setiap pilar.
  \item \textbf{Ukur kemampuan implementasi}: Soal menguji kemampuan mahasiswa untuk mengimplementasikan algoritma kriptografi dalam C.
  \item \textbf{Ukur kemampuan analisis kriptografi}: Soal meminta mahasiswa untuk menganalisis kelemahan algoritma dan melakukan serangan sederhana.
  \item \textbf{Ukur kemampuan sintesis}: Soal menuntut mahasiswa untuk merancang skema keamanan baru atau memodifikasi algoritma yang ada.
\end{enumerate}

Soal UTS yang efektif harus mengukur keseimbangan antara teori dan praktikum. Menurut sumber evaluasi, bobot soal yang terlalu sulit atau terlalu mudah dapat mengurangi motivasi dan pembelajaran yang efektif. Soal yang baik harus memiliki tingkat kesulitan yang sesuai dengan kemampuan mahasiswa dan memberikan kesempatan untuk menunjukkan pemahaman yang nyata.

Contoh soal UTS (style):
\begin{enumerate}
  \item Jelaskan tiga pilar keamanan informasi (confidentiality, integrity, authenticity) dan berikan contoh ancaman yang melanggar masing-masing.
  \item Hitung: (a) $17^{-1} \pmod{26}$, (b) $3^{10} \bmod 7$ menggunakan fast exponentiation.
  \item Enkripsi "CRYPTO" dengan Caesar $k=7$. Jelaskan mengapa Caesar cipher rentan terhadap frequency analysis.
  \item Jelaskan perbedaan struktur DES (Feistel) dan AES. Sebutkan empat operasi dalam satu round AES.
\end{enumerate}

Rubrik dan metode penilaian diuraikan pada section berikut.

\subsection{Soal UAS (Style)}

\begin{enumerate}
  \item Jelaskan langkah pembangkitan kunci RSA. Dengan $p=5$, $q=11$, $e=3$, hitung $n$, $\varphi(n)$, $d$, dan enkripsi $m=9$.
  
  \item Jelaskan peran fungsi hash dalam tanda tangan digital. Mengapa kita hash pesan sebelum menandatanganinya?
  
  \item Jelaskan protokol Diffie-Hellman. Mengapa Eve tidak dapat memperoleh kunci bersama meskipun melihat $A$ dan $B$?
  
  \item Implementasikan dalam C: fungsi \code{caesar\_encrypt} dengan parameter \code{char *text} dan \code{int shift}. Lampirkan kode lengkap.
\end{enumerate}

\section{Rubrik Penilaian Komprehensif}

Penilaian komprehensif memerlukan pendekatan yang sistematis dan objektif. Menurut sumber evaluasi, rubrik yang baik harus mencakup berbagai aspek:
\begin{itemize}
  \item \textbf{Kriteria penilaian}: Kriteria harus jelas, terukur, dan dapat diukur secara konsisten.
  \item \textbf{Feedback konstruktif}: Memberikan umpan balik yang mendalam dan konstruktif untuk perbaikan.
  \item \textbf{Authenticity}: Memastikan keaslian penilaian dan mencegah plagiarisme.
  \item \textbf{Fairness}: Menjamin proses penilaian yang adil dan transparan.
\end{itemize}

Rubrik penilaian yang baik membantu meningkatkan kualitas pembelajaran dan memastikan pencapaian yang objektif dan terukur. Menurut sumber evaluasi, implementasi rubrik yang konsisten dan berbasis bukti sangat penting untuk pengembangan program studi yang berkelanjutan. Pendekatan yang sistematis dan objektif memastikan bahwa evaluasi tidak hanya mengukur hasil belajar, tetapi juga mengembangkan kemampuan profesional dan etika akademis mahasiswa.

\subsection{Metode Penilaian}

Berbagai metode penilaian dapat digunakan untuk mengukur pencapaian:
\begin{itemize}
  \item \textbf{Observasi}: Pengamatan langsung proses pembelajaran dan partisipasi mahasiswa.
  \item \textbf{Portofolio}: Evaluasi karya mahasiswa untuk menunjukkan kemajuan praktikum.
  \item \textbf{Presentasi}: Presentasi hasil proyek atau penelitian.
  \item \textbf{Tes Tertulis}: Uji pemahaman dengan soal esai atau praktikum.
  \item \textbf{Peer Review}: Evaluasi kerja mahasiswa oleh teman sekelas.
  \item \textbf{Self-Assessment}: Refleksi diri dari pembelajaran.
\end{itemize}

Metode penilaian yang baik dirancang untuk memberikan feedback yang konstruktif dan mendukung pengembangan. Menurut sumber evaluasi, kombinasi berbagai metode ini memberikan gambaran yang komprehensif tentang kinerja dan pemahaman mahasiswa, serta memfasilitasi perbaikan berkelanjutan.

\textbf{Kriteria nilai}: A (85--100), B (70--84), C (60--69), D (50--59), E (<50).

\section{Tinjauan Pencapaian Kompetensi}

Gunakan checklist berikut untuk refleksi diri terhadap pencapaian CPMK:

\begin{checklist}
  \item Saya dapat menjelaskan prinsip dasar kriptografi dan keamanan informasi
  \item Saya dapat menerapkan aritmetika modular dan algoritma Euclid
  \item Saya dapat mengimplementasikan Caesar dan Vigenère dalam C
  \item Saya dapat menjelaskan DES, AES, dan mode operasi
  \item Saya dapat mengimplementasikan RSA dalam C (bilangan kecil)
  \item Saya dapat menjelaskan fungsi hash dan tanda tangan digital
  \item Saya dapat menganalisis protokol Diffie-Hellman
  \item Saya dapat melakukan frequency analysis dan Kasiski examination
  \item Saya dapat membuat dokumentasi dan pengujian kode C
\end{checklist}

Jika ada item yang belum tercapai, ulangi materi pada bab terkait dan kerjakan latihan tambahan.


% ============================================================
% AKTIVITAS PEMBELAJARAN
% ============================================================

\begin{aktivitas}
  \item \textbf{RNG}: Bandingkan output \texttt{rand()} dengan CSPRNG (\texttt{/dev/urandom} atau OpenSSL). Amati pola dengan visualisasi sederhana.

  \item \textbf{Praktikum C}: Implementasikan LCG dan uji periode untuk parameter berbeda.

  \item \textbf{Steganografi}: Implementasikan LSB embedding pada citra BMP/PNG. Sembunyikan pesan pendek, ekstrak, dan verifikasi.

  \item \textbf{Steganalis}: Diberikan citra asli dan citra dengan LSB stego, bandingkan histogram atau statistik LSB.
\end{aktivitas}

% ============================================================
% LATIHAN DAN REFLEKSI
% ============================================================

\begin{latihan}
  \item Bedakan PRNG dan TRNG. Kapan masing-masing digunakan?

  \item Mengapa LCG tidak aman untuk kriptografi? Apa prinsip Blum Blum Shub?

  \item Jelaskan uji keacakan NIST. Apakah lolos uji menjamin keamanan kriptografis?

  \item Jelaskan metode LSB untuk steganografi. Berapa kapasitas maksimum untuk citra 800×600 RGB?

  \item Bedakan steganografi dan watermarking.

  \item Implementasikan dalam C/Python: LSB embedding pada file citra.
\end{latihan}

% ============================================================
% ASESMEN
% ============================================================

\begin{asesmen}
\textbf{Instrumen Penilaian untuk Sub-CPMK 6.1, 7.1}

\textbf{A. Essay}

\begin{enumerate}
  \item Bedakan PRNG dan TRNG. Mengapa \texttt{rand()} tidak boleh dipakai untuk kunci?
  \item Jelaskan prinsip LSB steganografi. Apa kelemahannya terhadap steganalis?
\end{enumerate}

\textbf{B. Implementasi}

Implementasikan dalam C/Python: program yang menyembunyikan pesan teks dalam citra BMP/PNG dengan metode LSB, lalu mengekstrak kembali. Uji dengan citra dan pesan contoh.
\end{asesmen}

% ============================================================
% CHECKLIST KOMPETENSI
% ============================================================

\begin{checklist}
  \item Saya dapat membedakan PRNG dan TRNG
  \item Saya memahami LCG dan Blum Blum Shub
  \item Saya memahami pentingnya uji keacakan dan CSPRNG untuk kriptografi
  \item Saya dapat mengimplementasikan metode LSB untuk steganografi pada citra
  \item Saya memahami perbedaan steganografi dan watermarking
\end{checklist}

% ============================================================
% RANGKUMAN
% ============================================================

\begin{rangkuman}
Bab ini membahas pembangkit bilangan acak (TRNG vs PRNG, uji keacakan, Blum Blum Shub) dan steganografi (LSB, kapasitas, deteksi). RNG krusial untuk keamanan kriptografi; steganografi menyembunyikan informasi dalam media lain.

\textbf{Poin Kunci:}
\begin{itemize}
  \item TRNG: sumber fisik, tidak dapat diprediksi; PRNG: algoritma deterministik
  \item Uji keacakan: NIST SP 800-22, statistical tests
  \item Blum Blum Shub: CSPRNG yang dapat dibuktikan keamanannya
  \item LSB steganografi: embedding dalam bit tidak signifikan
\end{itemize}

\textbf{Kata Kunci}: \crypt{TRNG}, \crypt{PRNG}, \crypt{Blum Blum Shub}, \crypt{Steganografi}, \crypt{LSB}, \crypt{Keacakan}
\end{rangkuman}

% ============================================================
% Persiapan untuk Bab Berikutnya
% ============================================================

\begin{transition}
Setelah mempelajari kriptanalisis dan topik lanjutan, kita siap untuk memasuki \textbf{Bab X: Evaluasi dan Integrasi Kompetensi}.

Dalam Bab X, kita akan mengintegrasikan semua pengetahuan yang telah dipelajari:
\begin{itemize}
  \item \textbf{Integrasi Konsep}: Menghubungkan semua bab dalam sistem lengkap
  \item \textbf{Studi Kasus}: Menganalisis sistem keamanan nyata
  \item \textbf{Evaluasi Kompetensi}: Validasi pencapaian semua Sub-CPMK
  \item \textbf{Refleksi Pembelajaran}: Persiapan untuk aplikasi lanjutan
\end{itemize}

\textbf{Integrasi Semua Bab (II-IX)}:
\begin{itemize}
  \item \textbf{Bab II}: Layanan keamanan — fondasi kebutuhan keamanan
  \item \textbf{Bab III}: Matematika — alat untuk implementasi
  \item \textbf{Bab IV-V}: Cipher — teknik enkripsi klasik dan modern
  \item \textbf{Bab VI}: Kunci publik — solusi distribusi kunci
  \item \textbf{Bab VII}: Hash dan MAC — integrity dan authenticity
  \item \textbf{Bab VIII}: Protokol — implementasi praktis
  \item \textbf{Bab IX}: Kriptanalisis — perspektif keamanan
\end{itemize}

\textbf{Sistem Keamanan Lengkap}:
\begin{enumerate}
  \item \textbf{Identifikasi}: Sertifikat digital dan PKI
  \item \textbf{Enkripsi}: AES untuk data, RSA untuk kunci
  \item \textbf{Integrity}: SHA-256 dan HMAC
  \item \textbf{Authenticity}: Tanda tangan digital
  \item \textbf{Protokol}: TLS untuk komunikasi
  \item \textbf{Manajemen}: Siklus hidup kunci
\end{enumerate}

\textbf{Preview}: Di Bab X, kita akan melihat bagaimana semua komponen kriptografi bekerja sama dalam sistem keamanan nyata. Kita akan menganalisis studi kasus seperti HTTPS, email aman, dan sistem pembayaran. Kita juga akan mengevaluasi pemahaman kita tentang semua konsep dan mempersiapkan diri untuk aplikasi kriptografi dalam dunia nyata atau penelitian lebih lanjut.

Setelah Bab X, Anda akan memiliki pemahaman komprehensif tentang kriptografi — dari matematika dasar hingga implementasi praktis, dari desain cipher hingga analisis keamanan. Ini adalah fondasi yang kuat untuk karir di keamanan siber atau penelitian lanjutan dalam bidang kriptografi.
\end{transition}

\ifSubfilesClassLoaded{
  \renewcommand{\bibname}{Daftar Pustaka}
  \bibliographystyle{plain}
  \bibliography{../references}
}{}
\% end{document}




